% Adapted with attribution from the atomic-transport draft at the pinned repository snapshot.
\section{The short numerical construction}
\label{sec:numerical}

The following direct compiler is inherited from the successor-recursion
draft\cite{recursor}; it is included to make the present construction
self-contained. Atomic equality reasoning already gives transitivity
and numerical congruence. Define
\[
 W(n):=R(n,0;u,v.u),\qquad
 C(n,a,b):=W(R(n,a;u,v.Sb)),\qquad
 P(n):=C(n,0,W(n)).
\]
Computation and derived equality reasoning prove
\begin{align*}
 W0&=1,&W(Sn)&=n,\\
 C(0,a,b)&=a,&C(Sn,a,b)&=b,\\
 P0&=0,&P(Sn)&=n.
\end{align*}
The counter input $u$ of R is essential: it is not an alias for $v$.
Applying W proves that successor is injective.

Suppose fixed terms $g(\vec p)$ and $h(n,z,\vec p)$ are already
represented. Put
\[
 J_{g,h}(n,\vec p)=
 W\bigl(R(n,g(\vec p);u,v.S h(u,W(v),\vec p))\bigr).
 \tag{Compiler}\label{eq:compiler}
\]
Then
\[
 J_{g,h}(0,\vec p)=g(\vec p),\qquad
 J_{g,h}(Sn,\vec p)=h(n,J_{g,h}(n,\vec p),\vec p).
\]
At zero, R returns $Sg$ and W removes the successor. At a successor,
R returns $S h(n,J_{g,h}(n,\vec p),\vec p)$ and W again removes
the displayed successor. Both identities hold at arbitrary source
elements; there is no standard-numeral restriction or positivity premise.
Zero, successor, and projections are terms; composition is substitution.
External recursion on any fixed PR description therefore compiles it.
No scalar predecessor search, pairing construction, or equation induction
is needed for this compiler. Scalar iteration is its special case with
$h$ independent of its counter.

\section{Deriving the full intuitionistic logical interface}
\label{sec:logic}

This section adapts the inherited Choice-DI interface\cite{elimination,recursor}.
Every retained object formula is available by the Hyp rule, and
every compound equality transport means Theorem~\ref{thm:atomic-transport}.
Thus none of the following reasoning adds a primitive checker rule.

\subsection{Cases, equation induction, and congruence of the binder}

With an unused coordinate $x$, progress DI introduces $\D_{x,y}E(y)$
from $\Gamma,E(z)\vdash E(Sz)$ for a literal equation E. DE opens
its clause $\Q_yE(y)$; a supplied base proof $E(0)$ and QE at t give
$E(t)$. Hence contextual equation induction is admissible:
\[
 \frac{\Gamma\vdash E(0)\qquad\Gamma,E(z)\vdash E(Sz)}
      {\Gamma\vdash E(t)}. \tag{EI}\label{rule:EI}
\]
Case DI in the same construction, ignoring its base assumption in a
given successor proof, yields numerical cases for any object template:
\[
 \frac{\Gamma\vdash A(0)\qquad\Gamma\vdash A(Sz)}
      {\Gamma\vdash A(t)}. \tag{NC}\label{rule:NC}
\]
The usual freshness conditions are retained. Witness 1 may be used
for the unused coordinate. These derivations never need a proof whose
conclusion is an internal Q clause.

Uniform pointwise interderivability of matrices is preserved under D.
Eliminate the source D at a fresh witness, introduce the target D at
that witness by case DI, convert its base backwards, use the opened
source clause at $Sz$, and convert the result forwards. This is the
same constructive pattern as the transport theorem, now with arbitrary
uniform conversion derivations. Reverse the conversions for the converse.

\subsection{Numerical gates and formula selectors}

For a numerical control z define, externally on fixed formula syntax,
\[
 G_z(s=t):=(C(z,0,s)=C(z,0,t)),\qquad
 G_z(\D_{x,y}B):=\D_{x,y}G_zB.
\]
Freshen binders away from the control. Computation, reflexivity, and
binder congruence prove $\vdash G_0A$ and $G_{Sz}A\dashv\vdash A$.
The padding operation $\operatorname{Pad}(A)=\D_{u,y}G_yA$, with
unused u,y, is equivalent to A: case DI proves the forward direction;
DE, the true zero gate, and QE at 1 prove the reverse direction.

Define the selector on two equation leaves by
\[
 M_z(s=t,r=w):=(C(z,s,r)=C(z,t,w)).
\]
At aligned fresh binders extend it by
\begin{align*}
 M_z(\D_{x,y}A,\D_{x,y}B)&=\D_{x,y}M_z(A,B),\\
 M_z(E,\D_{x,y}B)&=\D_{x,y}M_z(G_yE,B),\\
 M_z(\D_{x,y}A,E)&=\D_{x,y}M_z(A,G_yE),
\end{align*}
where E is an equation. The maximum input binder depth decreases in
each recursive call. Padding and binder congruence give
\[
             M_0(A,B)\dashv\vdash A,\qquad
             M_{Sz}(A,B)\dashv\vdash B.
\]
These are finite syntactic operations, not object-language function symbols.

\subsection{Connectives and quantifiers}

Use the scalar iteration I compiled by \eqref{eq:compiler}, put
$Z(n)=C(n,1,0)$ and $p(n)=I(n,0;v.Z(v))$. The only parity facts
needed are the finite computations $p(1)=1$ and $p(2)=0$.
With fresh bound variables and unused u define
\begin{align*}
 A\Rightarrow_*B&=\D_{u,z}M_z(A,B),\\
 A\wedge_*B&=\D_{u,z}G_zM_{p(z)}(A,B),\\
 \exists_*x A(x)&=\D_{x,y}G_yA(x),\\
 A\vee_*B&=\exists_*x M_x(A,B),\\
 \forall_*x A(x)&=A(0)\wedge_*\D_{u,z}G_zA(z),\\
 \bot_*&=(0=1).
\end{align*}

Implication introduction converts its retained zero selector into A,
applies the supplied derivation of B, and converts B to the successor
selector; case DI closes. Elimination is DE followed by QE at 1.
Existential introduction uses its supplied witness and the successor
gate. Elimination opens D, supplies the theorem given by its zero gate,
and uses QE at 1. These are the usual eigenparameter conditions.

Disjunction introductions use control witnesses 0 and 1. For elimination,
turn the two branch derivations into $A\Rightarrow_*F$ and
$B\Rightarrow_*F$. NC gives $M_z(A,B)\Rightarrow_*F$ at any control;
eliminate the existential selector and apply the implication.
For conjunction introduction NC gives $M_v(A,B)$ at every v from the
two supplied proofs. Specialize to $p(Sz)$, add the successor gate,
and use case DI. Projections are QE at 1 and 2 after DE, using the two
displayed parity values.

NC now gives the coverage formula
\[
             (x=0)\vee_*\exists_*y\,(x=Sy).
\]
At zero use the left injection of reflexivity; at a successor use the
right injection with its evident witness. Coverage permits a case split
even if x occurs in retained assumptions: it adds an equality naming
the case, rather than illicitly substituting into the whole context.
Universal introduction obtains the zero and successor cases of a
fresh-parameter proof and forms the two conjuncts. Universal elimination
uses coverage, the first conjunct at zero, and the gated closure at a
successor. This avoids an eigenvariable assumption about the retained context.

Finally, under $0=1$, computation and atomic transport prove
\[
                   s=C(0,s,t)=C(1,s,t)=t
\]
for arbitrary s,t. Every equation follows; external recursion on object
syntax, using case DI at each D, proves every object formula. This
derives explosion for the designated absurdity. We have thus derived
the full intuitionistic natural-deduction rules, with actual equality,
without postulating excluded middle for quantified formulas.

\section{Full arithmetic and both interpretation composites}
\label{sec:interpretations}

\subsection{An explicit ordinary realization}

Let $T_R$ be ordinary intuitionistic first-order logic with equality
in the staged R signature, the equations R0 and RS, contextual
literal-equation induction EI, and
\[
 \forall x\,(x=0\vee\exists y\,x=Sy),\qquad (0=1)\to\bot.
 \tag{Boundary}
\]
For definiteness, ordinary logic here has hypothesis, implication
introduction and modus ponens, conjunction introduction and its two
projections, the two disjunction introductions and case elimination,
universal introduction and specialization, existential introduction and
fresh-witness elimination, and explosion. Equality has reflexivity and
substitution. The usual eigenvariable restrictions are imposed. These
ordinary rules specify the comparison theory, not additional primitives
of $\Acal$.

Realize its object and clause syntax by
\[
 (s=t)^\circ=(s=t),\quad
 (\D_{x,y}B)^\circ=\exists x(B^\circ(x,0)\to\forall yB^\circ(x,y)),
\]
\[
                 (\Q_yB)^\circ=B^\circ(0)\to\forall yB^\circ(y).
\]
The interpretation domain is $x=x$ and equality is actual equality.
Case DI is sound by coverage. Progress DI is sound by equation
induction under its temporarily assumed base. All other native rules
have their direct ordinary meanings. Conversely Section~\ref{sec:logic}
derives the ordinary rules, coverage, and EI. Separation translates to
the identity implication $(0=1)\Rightarrow_*(0=1)$.

The explicit falsity bridge matters. The positive numerical equations
alone admit a singleton structure, so declaring them an ordinary
first-order theory while silently omitting separation would be a mistake.

\subsection{Actual equality reflection}

For clarity, the remaining arithmetic bridge is given internally,
not merely by agreement on standard numerals. Every fixed PR description
is already compiled by Section~\ref{sec:numerical}. In particular use
\[
 x+0=x,\quad x+Sy=S(x+y),\qquad
 d(x,0)=x,\quad d(x,Sy)=P(d(x,y)).
\]
Let $\operatorname{sg}(x)=C(x,0,1)$. Coverage and the computation
equations show that sg takes only 0,1, that
$\operatorname{sg}(x)=0\leftrightarrow x=0$, and $C(x,a,a)=a$.
Applying sg and separation gives $Sx\ne0$. Coverage and W give
$W(x)=0\leftrightarrow x=1$.

Literal-equation induction proves, in order,
\[
 0+y=y,\qquad (Sx)+y=S(x+y),\qquad
 d(Sx,Sy)=d(x,y),\qquad d(Sa,a)=1.
\]
For the first two induct on y with x fixed. For the third induct on y:
the base is $P(Sx)=x$, and its step applies P to the induction equation.
For the fourth induct on a and use the third identity in the step.

One more equation induction, on y with x fixed, gives
\[
        C(\operatorname{sg}(d(x,y)),x,y+d(x,y))=x. \tag{Residual}
\]
The base is $C(\operatorname{sg}(x),x,x)=x$. In the step set
$r=d(x,y)$ and $r'=Pr=d(x,Sy)$. Coverage splits $r'$ into zero
or successor. At zero the selector already returns x. In the successor
case, $r'\ne0$ forces r to be a successor, and $P(Sq)=q$ gives
$r=Sr'$. The induction equation says $y+r=x$. The new second branch
is $(Sy)+r'=y+Sr'=y+r=x$. This completes the induction without
using a decidable actual-equality predicate in its premises.

Consequently $d(Sa,b)=1\leftrightarrow a=b$: forwards, Residual yields
$b+1=Sa$, hence $Sb=Sa$, and W cancels the successors; backwards
use $d(Sa,a)=1$. Thus
\[
                 e(a,b)=\operatorname{sg}(W(d(Sa,b)))
\]
is Boolean and satisfies $e(a,b)=0\leftrightarrow a=b$ for arbitrary
elements of every source model.

For a fixed quantifier-free formula A, define its zero-true
characteristic by
\begin{align*}
 b_{s=t}&=e(s,t),&b_\bot&=1,\\
 b_{A\wedge B}&=C(b_A,b_B,1),&
 b_{A\vee B}&=C(b_A,0,b_B),&
 b_{A\to B}&=C(b_A,b_B,0).
\end{align*}
An external induction on A, using the proved Boolean ranges, gives
$b_A\in\{0,1\}$ and $b_A=0\leftrightarrow A$ intuitionistically.
The implication case uses that $b_A=1$ proves $\neg A$; it does not
assert excluded middle for arbitrary quantified formulas. Applying EI
to $b_A(n)=0$ recovers every contextual quantifier-free induction instance.

\subsection{The two numerical interpretations}

Take iPRA in its full first-order language of fixed PR descriptions,
with defining equations, constructor separation, and quantifier-free
induction. Its atomic-induction formulation is equivalent by the
characteristic argument above; this is the usual first-order iPRA
convention\cite{salehi}. Interpret its descriptions in $T_R$ by
\eqref{eq:compiler}, leaving domain, equality, and ordinary logic unchanged.
All defining equations, successor axioms, and induction instances have
just been verified internally.

Conversely, each fixed R body is interpreted in iPRA by ordinary
primitive recursion with base $Sa$ and step $h(u,v,\vec p)$. The
abbreviation $S(a)=R(0,a;u,v.0)$ then has ordinary successor value.
R0, RS, the boundary sentences, and EI are iPRA theorems.

For the stronger bi-interpretability claim, first anchor the successor
abbreviation. Its iPRA interpretation is the zero instance of the
recursor with zero update, so its defining base clause gives ordinary
$Sa$. Compiling that interpreted successor back gives the original
$S(a)$ by the compiler's base computation. This includes the expansion
$r_0(0,a)$ used to name successor, and does not presuppose uniqueness
of all R symbols. Thus the same successor occurs in both copies of each
recursion before the following staged uniqueness argument begins.

Now fix independent initializer and parameter variables. Two solutions
of the same fixed R recursion
are internally equal by induction on their value equation: bases agree,
and their steps agree by equality congruence and the induction equation.
Induct externally on the finite staged description. Its double
translation satisfies the same equations once earlier subdescriptions
have been identified, so uniqueness identifies it with the original.
The identical argument works for ordinary PR descriptions. Only after
these identities have been proved for independent parameters do we
substitute counter-dependent actual arguments. This also establishes
coherence when substitution changes the name of a staged body.

Both numerical interpretation domains are the entire original domain;
both internally definable isomorphisms have graph $x=y$. Nothing is
restricted to the standard cut or quotiented by a numerical observer.

\subsection{The logical round trips}

The ordinary meanings of gates and selectors give the advertised
meanings of all six macros in Section~\ref{sec:logic}, with coverage
in the universal case and separation for falsity. For the only primitive
object binder the reverse composite is explicitly
\[
 \D_{x,y}B\dashv\vdash
 \exists_*x\bigl(B(x,0)\Rightarrow_*\forall_*yB(x,y)\bigr).
\]
Forwards, open D at fresh a, use its available clause under the assumed
base to prove each $B(a,y)$, form the conditional universal, and introduce
the existential at a. Backwards, eliminate the existential at fresh a,
introduce the original D by case DI, and use its retained base, the
conditional, and universal specialization to obtain the successor case.
Neither direction requires Q to be an object goal. External induction
on formulas establishes both logical composites.

\begin{theorem}[Whole-domain bi-interpretation]
The ordinary realization of $\Acal$ is bi-interpretable with full
first-order iPRA. The two numerical isomorphisms are internally given
by $x=y$, and both logical formula composites are provably equivalent
to their originals. Effective finite proof translations preserve the
designated contradiction, hence the two theories are equiconsistent.
\end{theorem}
This is an external metatheorem. No claim is made that either theory
proves its own consistency or the uniform correctness of an arbitrary
primitive-recursive evaluator.

